\subsection{Basics}

riemann integral =(.
``im kind of biased in saying this, but im also correct, so". fire.

\begin{defn}
	for an int'vl $I\subseteq\bR$ w endpts $a<b$, define $\ell(I)=b-a$.
	for $A\subseteq\bR$, its \textbf{outer measure} $m^{*}(A)$ is:
	\begin{equation*}
		m^{*}(A) \ := \ \inf\set{\sum_{k=1}^{\infty}\ell(I_{k}): A\subseteq\bigcup_{k=1}^{\infty}I_{k}, \  I_{k} \trm{ open bdd int'vls}}
	\end{equation*}
\end{defn}
note $m^{*}(A)\in[0,\infty]$, i.e. $\bR_{\geq0}$ but can also be infty.

\begin{prop} \vspace{-0.2in}
	\begin{enumerate}[i)]
		\item if $A\subseteq B$, $m^{*}(A)\leq m^{*}(B)$
		\item if $E$ ctbl, $m^{*}(E)=0$
		\item if $I$ int'vl, then $m^{*}(I)=\ell(I)$
	\end{enumerate}
\end{prop} \

\begin{pf} \vspace{-0.2in}
	\begin{enumerate}[i)]
		\item any cvr of $B$ cvrs $A$.
		\item if $E=\set{x_{i}}$, let $I_{k}=(x_{k}-\ep/2^{k+1},x_{k}+\ep/2^{k+1})$. then:
			\begin{equation*}
				m^{*}(E)\leq\sum\ell(I_{k})=\sum\ep/2^{k}=\ep
			\end{equation*}
		\item sps $J=[a,b]$ (renamed). sps $I_{k}$ open int'vls cvring $J$.
			$J$ cpt, so take finitely many; also assume no $I_{j}\subseteq I_{k}$.
			if $I_{k}=(a_{k},b_{k})$ and $\set{a_{i}},\set{b_{i}}$ increasing, then:
			\begin{equation*}
				\sum_{j=1}^{k}\ell(I_{j}) \ = \ \sum_{j=1}^{k}b_{j}-a_{j}
				\ = \ b_{k}-a_{1}+\sum b_{j}-a_{j+1} \ > \ b_{k}-a_{1} \ \geq \ b-a
			\end{equation*}
			so $\fa$ cvrs $\set{I_{j}}$, $\sum\ell(I_{j})\geq b-a$ implies $m^{*}(J)\geq b-a$.
			OTOH, if $J_{\ep}=(a-\ep,b+\ep)$, then $m^{*}(J)\geq\ell(J_{\ep})\sto b-a$.
			[i love aboos of notation]

			now, if $J=(a,b)$, let $I^{-}=[a+\ep,b-\ep]$ and $I^{+}=[a-\ep,b+\ep]$.
			clearly, $I^{-}\subseteq J\subseteq I^{+}$,
			so $m^{*}(I^{-})\leq m^{*}(J)\leq m^{*}(I^{+})$, and it follows.
	\end{enumerate}
\end{pf} \

\begin{lm}
	if $E\subseteq\bR$ and $y\in\bR$, then $m^{*}(E+y)=m^{*}(E)$.
\end{lm}
translate ur cvrs.

\begin{lm}
	if $E\subseteq\bigcup_{k=1}^{\infty}E_{k}$, then $m^{*}(E)\leq\sum m^{*}(E_{k})$.
\end{lm} \

\begin{pf}
	for all $k$, let $\set{I_{j}^{(k)}}$ be an int'vl cvr of $E_{k}$ s.t.:
	\begin{equation*}
		m^{*}(E_{k}) \ \leq \ \sum_{j=1}^{\infty}\ell(I_{j}^{(k)}) \ \leq \ m^{*}(E_{k})+\frac{\ep}{2k}
	\end{equation*}
	since $\set{I_{j}^{(k)}}_{j,k}$ an int'vl cvr of $E$, then:
	\begin{equation*}
		m^{*}(E) \ \leq \ \sum_{k=1}^{\infty}\sum_{j=1}^{\infty}\ell(I_{j}^{(k)})
		\ \leq \ \sum_{k}m^{*}(E_{k})+\frac{\ep}{2k}
		\ \leq \ \sum_{k}m^{*}(E_{k}) \ + \ \ep
	\end{equation*}
\end{pf} \

\begin{defn}
	a set $E\subseteq\bR$ is \textbf{measurable} if $\fa A\subseteq\bR$:
	\begin{equation*}
		m^{*}(A) \ = \ m^{*}(A\cap E) + m^{*}(A\cap E^{c})
	\end{equation*}
	we denote by $\cl{E}$ the set of measurable sets.
\end{defn}
note by the union property, we always have 
$m^{*}(A) \leq m^{*}(A\cap E) + m^{*}(A\cap E^{c})$.
thus, to show $E$ measurable, we only need to show:
\begin{equation*}
	m^{*}(A) \ \geq \ m^{*}(A\cap E) + m^{*}(A\cap E^{c})
\end{equation*}
note also, non-measurable sets exist; that is, $\ex A,B$ disjoint s.t.:
\begin{equation*}
	m^{*}(A\cup B) \ \lneq \ m^{*}(A) + m^{*}(B)
\end{equation*}

\begin{lm}
	if $m^{*}(E)=0$, then $E\in\cl{E}$.
\end{lm} \

\begin{pf}
	note
	$m^{*}(A\cap E)+m^{*}(A\cap E^{c}) \ \leq \ m^{*}(E)+m^{*}(A) \ = \ m^{*}(A)$.
\end{pf} \

\begin{lm}
	if $E_{1},E_{2}\in\cl{E}$, then $E_{1}\cup E_{2}\in\cl{E}$.
\end{lm} \

\begin{pf}
	sps $A\subseteq\bR$. then:
	\begin{align*}
		m^{*}(A)
		& \ = \ m^{*}(A\cap E_{1}) + m^{*}(A\cap E_{1}^{c}) \\
		& \ = \ m^{*}(A\cap E_{1}) + m^{*}(A\cap E_{1}^{c}\cap E_{2}) + m^{*}(A\cap E_{1}^{c}\cap E_{2}^{c}) \\
		& \ \geq \ m^{*}(A\cap(E_{1}\cup E_{2})) + m^{*}(A\cap(E_{1}\cup E_{2})^{c})
	\end{align*}
	since $A\cap(E_{1}\cup E_{2}) \ = \ (A\cap E_{1})\cup(A\cap E_{1}^{c}\cap E_{2})$.
\end{pf} \

\begin{defn}
	a coll'n $\cl{A}\subseteq\cl{P}(\bR)$ is an \textbf{algebra} if:
	\begin{itemize}
		\item $\bR\in\cl{A}$
		\item $E\in\cl{A} \ \implies \ E^{c}\in\cl{A}$
		\item $E_{1}\cup E_{2}\in\cl{A} \ \implies \ E_{1}\cup E_{2}\in\cl{A}$
	\end{itemize}
\end{defn} \

\begin{crll}
	$\cl{E}$ is an alg
\end{crll} \

\begin{prop}
	if $C_{1},C_{2}\in\cl{E}$ disj, then $m^{*}(C_{1}\cup C_{2})=m^{*}(C_{1})+m^{*}(C_{2})$.
\end{prop} \

\begin{pf}
	note $A\cap E_{1} \ = \ A\cap E_{1}\cap E_{2}+A\cap E_{1}\cap E_{2}^{c}$, so:
	\begin{align*}
		m^{*}(A\cap(E_{1}\cup E_{2}))
		& \ = \ m^{*}(A\cap(E_{1}\cup E_{2})\cap E_{2})+m^{*}(A\cap(E_{1}\cup E_{2})\cap E_{2}^{c}) \\
		& \ = \ m^{*}(A\cap E_{1}) + m^{*}(A\cap E_{2})
	\end{align*}
\end{pf} \

\begin{prop}
	$\cl{E}$ closed under ctbl unions.
\end{prop} \

\begin{pf}
	let $A_{j}\in\cl{E}$. denote:
	\begin{equation*}
		E_{1} \ = \ A_{1} \qquad E_{j} \ = \ A_{j}\sm \bigcup_{k=1}^{j-1}A_{k}
	\end{equation*}
	then $E_{j}$ disj and $E=\bigcup E_{j}$.
	let $F_{n}=\bigcup_{j=1}^{n}E_{j}$. then:
	\begin{align*}
		m^{*}(A)
		& \ = \ m^{*}(A\cap F_{n}) + m^{*}(A\cap F_{n}^{c}) \\
		& \ \geq \ m^{*}(A\cap F_{n}) + m^{*}(A\cap E^{c}) \\
		& \ = \ \sum_{j=1}^{n}m^{*}(A\cap E_{j}) + m^{*}(A\cap E^{c})
	\end{align*}
	taking $n\sto\infty$:
	\begin{equation*}
		m^{*}(A) \ \geq \ \sum_{j=1}^{\infty}m^{*}(A\cap E_{j})+m^{*}(A\cap E^{c})
		\ \geq \ m^{*}(A\cap E) + m^{*}(A\cap E^{c})
	\end{equation*}
\end{pf}

\begin{defn}
	a $\sigma$-algebra is an alg closed under ctbl unions.
\end{defn}
thus, $\cl{E}$ is a $\sigma$-alg.
note also that $\cl{E}$ closed under ctbl intersections by taking complements.

\begin{defn}
	we call $\cl{E}$ the set of \textbf{Lebesgue measurable} sets, and denote $m:=m^{*}\rvert_{\cl{E}}$
	as the \textbf{Lebesgue measure}.
\end{defn} \

\begin{lm}
	$(a,\infty)\in\cl{E}$.
\end{lm} \

\begin{pf}
	let $A\subseteq\bR$; first sps $a\notin A$.

	sps $\set{I_{j}}$ cvrs $A$.
	let $J_{k}=I_{k}\cap(-\infty,a),K_{k}=I_{k}\cap(a,\infty)$.
	note:
	\begin{equation*}
		\sum_{j}\ell(I_{j}) \ = \ \sum_{j}\ell(J_{j})+\sum_{j}\ell(K_{j})
	\end{equation*}
	we also have:
	\begin{equation*}
		A\cap(-\infty,a) \ \subseteq \ \bigcup_{j}J_{j} \qquad
		A\cap(a,\infty) \ \subseteq \ \bigcup_{j}K_{j}
	\end{equation*}
	therefore, we have that:
	\begin{equation*}
		m^{*}(A\cap(-\infty,a))+m^{*}(A\cap(a,\infty))
		\ \leq \ \sum_{j}\ell(J_{k})+\sum_{j}\ell(K_{j}) \ = \ \sum_{j}\ell(I_{j})
	\end{equation*}
	and it follows that $m^{*}(A\cap(-\infty,a))+m^{*}(A\cap(a,\infty)) \leq m^{*}(A)$.
\end{pf}

\newpage
\lecdate{Lec 16 - Mar 10 (Week 9)}
(missed)

\begin{crll}
	all open, half-open, closed int'vls are in $\cl{E}$
\end{crll} \

\begin{pf}
	note:
	\begin{equation*}
		[a,\infty) \ = \ \bigcap_{n\geq1}(a-n^{-1},\infty) \ \in \ \cl{E}
	\end{equation*}
	thus, any semi-infinite int'vl in $\cl{E}$, so
	$(a,b)=(a,\infty)\cap(-\infty,b)$ in $\cl{E}$.
	other int'vls follow similarly.
\end{pf} \

\begin{crll}
	any open set is in $\cl{E}$
\end{crll} \

\begin{thm}[title=Regularity]
	let $E\in\cl{E}$, $\ep>0$.
	then there are open and closed $O,F$ resp., s.t.:
	\begin{equation*}
		m(O\sm E)<\ep \qquad m(E\sm F)<\ep
	\end{equation*}
\end{thm} \

\begin{pf}
	we start with outer reg. assume $m(E)<\infty$.
	then there are int'vls $I_{k}$ s.t.:
	\begin{equation*}
		\sum_{k}m(I_{k}) \ < \ m(E)+\ep
	\end{equation*}
	then $O=\bigcup_{k}I_{k}$ works since:
	\begin{equation*}
		m(O\sm E) \ = \ m(O) - m(E) \ \leq \ \sum_{k}m(I_{k}) - m(E) \ < \ \ep
	\end{equation*}
	now sps $m(E)=\infty$. let $E_{k}=E\cap[-k,k]$. then $m(E_{k})\leq2k$.
	choose sets $O_{k}\supseteq E_{k}$ s.t. $m(O_{k}\sm E_{k})<\frac{\ep}{2^{k}}$.
	let $O=\bigcup_{k}O_{k}$. then:
	\begin{equation*}
		m(O\sm E) \ \leq \ \sum_{k}m(O_{k}\sm E_{k}) \ < \ \ep
	\end{equation*}
	for inner reg., choose $O\supseteq E^{c}$ s.t.
	$m(O\sm E^{c})<\ep$ and let $F=O^{c}$. then:
	\begin{equation*}
		m(E\sm F) \ = \ m(O\sm E^{c}) \ < \ \ep
	\end{equation*}
	since $E\sm F \ = \ E\cap F^{c} \ = \ E\cap O \ = \ O\sm E^{c}$.
\end{pf} \

\begin{crll}
	for any $E\in\cl{E}$, there are open $O_{n}\supseteq E$ and closed $F_{n}\subseteq E$ s.t.:
	\begin{equation*}
		m\left(\bigcap_{n}O_{n}\sm E\right) = 0 \qquad
		m\left(E\sm\bigcup_{n}F_{n}\right) = 0
	\end{equation*}
\end{crll}
this is another way of saying that we measure a set by 
an open outer approx or closed inner approx.
here, we want the measure of the difference to go to 0.

\lecdate{Lec 17 - Mar 17 (Week 10)}
midterm last thurs

\begin{thm}\vspace{-0.2in}
	\begin{enumerate}[i)]
		\item if $A_{n}\subseteq A_{n+1}\in\cl{E}$ then
			$m\left(\bigcup A_{n}\right)=\lim_{n}m(A_{n})$
		\item if $m(B_{1})<\infty$ and $B_{n}\supseteq B_{n+1}$ then
			$m\left(\bigcap B_{n}\right)=\lim_{n}m(B_{n})$
	\end{enumerate}
\end{thm} \

\begin{pf} \vspace{-0.2in}
	\begin{enumerate}[i)]
		\item let $C_{1}=A_{1}$ and $C_{n}=A_{n}\sm A_{n-1}=A_{n}\cap(A_{n-1})^{c}\in\cl{E}$.
			then they are disj, so:
			\begin{gather*}
				m\left(\bigcup_{n=1}^{\infty}A_{n}\right)
				\ = \ m\left(\bigcup_{n=1}^{\infty}C_{n}\right)
				\ = \ \sum_{n=1}^{\infty}m(C_{n}) \\
				\ = \ \lim_{n\sto\infty}\sum_{k=1}^{n}m(C_{k})
				\ = \ \lim_{n\sto\infty}m\left(\bigcup_{k=1}^{n}C_{k}\right)
				\ = \ \lim_{n\sto\infty}m(A_{n})
			\end{gather*}
		\item let $A_{n}=B_{1}\sm B_{n}$.
			since $B_{n}$ is decreasing, $A_{n}$ is increasing, so
			\begin{equation*}
				m\left(\bigcup_{n=1}^{\infty}A_{n}\right)=\lim_{n\sto\infty}m(A_{n})
			\end{equation*}
			note:
			\begin{equation*}
				m(A_{n}) \ = \ m(B_{1}\sm B_{n}) \ = \ m(B_{1})-m(B_{n})
			\end{equation*}
			OTOH:
			\begin{align*}
				m\left(\bigcup_{n=1}^{\infty}A_{n}\right)
				& \ = \ m\left(\bigcup_{n=1}^{\infty}B_{1}\sm B_{n}\right) \\
				& \ = \ m\left(B_{1}\sm \bigcap_{n=1}^{\infty}B_{n}\right) \\
				& \ = \ m(B_{1})-m\left(\bigcap_{n=1}^{\infty}B_{n}\right) \\
				& \ = \ \lim_{n\sto\infty}m(B_{1})-m(B_{n})
			\end{align*}
	\end{enumerate}
\end{pf}

recall the defn of a $\sigma$-alg.
given a coll'n of subsets of $\bR$ the $\sigma$-alg \textit{generated by} $A$ is
the intersection of all $\sigma$-algs containing $A$, denoted $\sigma(A)$.

\begin{defn}
	a \textbf{measure space} $(\Omega,\cl{F})$ is a set $\Omega$ and a $\sigma$-alg $\cl{F}$ on $\Omega$.
	given two measure spaces $(\Omega_{1},\cl{F}_{1}),(\Omega_{2},\cl{F}_{2})$,
	we say $f:\Omega_{1}\gto\Omega_{2}$ is \textbf{measurable} if $\fa E\in\cl{F}_{2}$,
	we have that $f^{-1}(E)\in\cl{F}_{1}$.
\end{defn}
note the similarity to (topological) continuity,
not that the similarity is particularly important though.

\begin{lm}
	if $(\Omega_{1},\cl{F}_{1}),(\Omega_{2},\cl{F}_{2})$ are measure spaces
	s.t. $\cl{F}_{2}=\sigma(A)$ for some coll'n $A\subseteq\cl{P}(\Omega_{2})$,
	then $f:\Omega_{1}\gto\Omega_{2}$ is measurable if:
	\begin{equation*}
		f^{-1}(E) \ \in \ \cl{F}_{1}
	\end{equation*}
	for all $E\in A$.
\end{lm} \

\begin{pf}
	suffices to check that
	\begin{equation*}
		S \ := \ \set{E:f^{-1}(E)\in\cl{F}_{1}}
	\end{equation*}
	is a $\sigma$-alg containing $A$.
	indeed, its a $\sigma$-alg essentially bc $\cl{F}_{1}$ is.
	\begin{equation*}
		E_{n}\in S \ \implies \ f^{-1}\left(\bigcup_{n=1}^{\infty}E_{n}\right)
		= \bigcup_{n=1}^{\infty}f^{-1}(E_{n}) \in \cl{F}_{1}
		\ \implies \ \bigcup_{n=1}^{\infty}E_{n}\in S
	\end{equation*}
	other axioms are similar.
\end{pf}
so why do we care abt this lemma?

\begin{defn}
	the \textbf{Borel} $\bm{\sigma}$\textbf{-alg} is given as
	$\cl{B}(\bR) \ := \ \sigma(\trm{open subsets of } \bR)$.
\end{defn}
fact: $\cl{B}(\bR)\subsetneq\cl{E}\subsetneq2^{\bR}$.
the inclusions being strict is tedious and unnecessary to show,
and thus skipped.

\begin{defn}
	we say a fn $f:\bR\gto\bR$ is \textbf{Lebesgue measurable} if it is measurable as
	a fn $(\bR,\cl{E})\gto(\bR,\cl{B}(\bR))$.
	we say it is \textbf{Borel measurable} if it is measurable as
	a fn $(\bR,\cl{B}(\bR))\gto(\bR,\cl{B}(\bR))$.
\end{defn} \

\begin{lm}
	a fn $f:\bR\gto\bR$ is lebesgue (borel) meas if:
	\begin{equation*}
		f^{-1}\left((c,\infty)\right)\in\cl{E} \qquad (\trm{or } \cl{B}(\bR))
	\end{equation*}
	for any $c\in\bQ$.
\end{lm} \

\begin{pf}
	suffices to check $\set{\trm{open sets}}\subseteq\sigma(\set{(c,\infty)})=:\cl{F}$
	for $c\in\bQ$.
	\begin{enumerate}[i)]
		\item check $(a,b)\in\cl{F}$.
			indeed, $(a,b)=(a,\infty)\sm[b,\infty)$, but
			$[b,\infty)=\bigcap(q,\infty)$ for $q\in\bQ\cap(-\infty,b)$.
		\item any open set in $\bR$ is a ctbl union of open intvls w/
			rational endpts.
	\end{enumerate}
\end{pf} \

\begin{crll}
	any cts or monotonic $f:\bR\gto\bR$ is borel measurable.
\end{crll} \

\begin{prop}
	sps $f,g:\bR\gto\bR$ are lebesgue (borel) meas.
	the following are as well:
	\begin{enumerate}[i)]
		\item $\abs{f},\alpha f+\beta$ for real $\alpha,\beta$
		\item $(f)_{+},(f)_{-}$ where $(x)_{*}$ is the $+$ve or $-$ve parts of $x$
	\end{enumerate}
\end{prop} \

\begin{pf}
	these are compositions of $f$ w/ a cts fn, hence borel meas.
	easy to see that the composition of a BM fn w/ a LM fn is LM.
\end{pf}

\lecdate{Lec 18 - Mar 19 (Week 10)}

\begin{prop}
	in addition to the above, the following are also lebesgue (borel) meas.:
	\begin{enumerate}
		\item[iii)] $f+g,f-g,fg$
		\item[iv)] $\max(f,g),\min(f,g)$
	\end{enumerate}
\end{prop} \

\begin{pf} \vspace{-0.2in}
	\begin{enumerate}
		\item[iii)] we show $f+g$ is lebesgue.
			suffices to check $\set{x:f(x)+g(x)>c}\in\cl{E}$ $\fa c\in\bR$.
			we claim that:
			\begin{equation*}
				\set{x:f(x)+g(x)>c}\in\cl{E}
				\ = \ \bigcup_{q\in\bQ}\set{x:f(x)>c-q}\cap\set{x:g(x)>q}
			\end{equation*}
			enough to check $(\subseteq)$.
			indeed, sps $x$ is s.t. $f(x)+g(x)>c$.
			then $\ex q\in\bQ$ with $g(x)>q>c-f(x)$, with
			$f(x)>c-q$ and $g(x)>q$.

			$f-g$ is similar, so we show $fg$. note:
			\begin{equation*}
				fg \ = \ \frac{1}{4}\left((f+g)^{2}-(f-g)^{2}\right)
			\end{equation*}
			wts if $f$ meas, then $f^{2}$ is,
			but this is just a composition of meas. fns.
		\item [iv)] follows from above, as:
			\begin{equation*}
				\max(f,g) \ = \ \frac{\abs{f-g}+(f+g)}{2} \qquad
				\min(f,g) \ = \ \frac{(f+g)-\abs{f-g}}{2}
			\end{equation*}
	\end{enumerate}
\end{pf}

it is sometimes useful to consider fns which take values $\pm \infty$.
\begin{defn}
	let $\bar{\bR}=\bR\cup\set{\pm \infty}$ denote the extended real line.
	we define:
	\begin{equation*}
		\cl{B}(\bar{\bR}) \ := \
		\set{E\subseteq\bar{\bR}:E=A\cup B,A\in\cl{B}(\bR),B\in\set{\eset,\set{\infty},\set{-\infty},\set{\pm \infty}}}
	\end{equation*}
\end{defn} \

\begin{exr} \vspace{-0.2in}
	\begin{enumerate}[i)]
		\item $\cl{B}(\bar{\bR})$ is a $\sigma$-alg
		\item $f:\bR\gto\bar{\bR}$ is LM (BM) iff
			$\set{x:f(x)>c}\in\cl{E}$ or $\cl{B}(\bR)$ $\fa c\in\bQ$
		\item if $f:\bR\gto\bR$ meas., then $\iota\circ f:\bR\gto\bar{\bR}$ meas.
	\end{enumerate}
\end{exr}
any meas. fn defined on the ext. reals is (unsurprisingly) also called LM/BM.
the difference only rly matters when it comes to operations on fns (whats $\infty-\infty$?).

given $E\in\cl{E}$, $f:E\gto\bar{\bR}$, tfae:
\begin{itemize}
	\item $f$ meas
	\item $f^{-1}(A)\in\cl{E}$ or $\cl{B}(\bR)$ $\fa A\in\cl{B}(\bar{\bR})$
	\item the fn $\tilde{f} = \chi_{E}f$ is meas.
\end{itemize}
equivalence of last one is an exercise.

\subsection{Limits of measurable functions}

\begin{prop}
	let $f_{n}:E\gto\bar{\bR}$ a seq of ptwise incresing meas. fns for some
	$E$ meas.
	then $\lim_{n\sto\infty}f_{n}(x)=f(x)$ is meas.
\end{prop} \

\begin{pf}
	for all $c\in\bR$, since $f_{n}$ is ptwise increasing, we have:
	\begin{equation*}
		\set{x:f(x)>c} \ = \ \bigcap_{n=1}^{\infty}\set{x:f_{n}(x)>c}
	\end{equation*}
\end{pf} \

\begin{thm}
	sps $E$ meas., $f_{n}:E\gto\bar{\bR}$ seq of meas. fns.
	the following are also meas.:
	\begin{enumerate}[i)]
		\item $x\mto\sup_{n}f_{n}(x)$
		\item $x\mto\inf_{n}f_{n}(x)$
		\item $x\mto\limsup_{n}f_{n}(x)$
		\item $x\mto\liminf_{n}f_{n}(x)$
		\item $x\mto\lim_{n}f_{n}(x)$, provided this exists $\fa x$
	\end{enumerate}
\end{thm} \

\begin{pf}\vspace{-0.2in}
	\begin{enumerate}[i)]
		\item $\fa c\in\bR$:
			\begin{equation*}
				\set{x:\sup_{n}f_{n}(x)>c} \ = \ \bigcup_{n=1}^{\infty}\set{x:f_{n}(x)>c}
			\end{equation*}
		\item similar as above
		\item recall $\limsup_{n}f_{n}(x)=\lim_{n}\set{\sup_{m\geq n}f_{m}(x)}$.
			this is a decreasing seq of meas. fns, so by i), we're done
		\item similar as above
		\item if the lim exists, its equal to both limsup and liminf
	\end{enumerate}
\end{pf}

we know what a.e. is (uses outer meas. specifically, since we dk
a priori the set is meas)

\begin{lm}
	sps $f:\bR\gto\bar{\bR}$ LM and $g:\bR\gto\bar{\bR}$ is s.t. $f(x)=g(x)$ ae.
	then $g$ meas.
\end{lm} \

\begin{pf}
	let $E=\set{x:f(x)\neq g(x)}$.
	for $A\in\cl{B}(\bar{\bR})$, we have:
	\begin{equation*}
		g^{-1}(A) \ = \ (f^{-1}(A)\cap E^{c}) \cup F
	\end{equation*}
	where $F=g^{-1}(A)\cap E$.
	then $m^{*}(F)\leq m^{*}(E)=0$, so $F\in\cl{E}$.
	by assumption, $f^{-1}(A)\in\cl{E}$ and $E^{c}\in\cl{E}$,
	so $g^{-1}(A)\in\cl{E}$.
\end{pf} \

\begin{crll}
	sps $f_{n}:\bR\gto\bar{\bR}$ meas s.t. $\lim_{n}f_{n}(x)$ exists ae.
	let $f(x)=\lim_{n}f_{n}(x)$ when it exists, and arbitrarily otw.
	then $f$ LM.
\end{crll} \

\begin{pf}
	let $E=\set{x:\lim_{n}f_{n}(x) \trm{ exists}}$.
	note $E\in\cl{E}$ since $m^{*}(E^{c})=0$.
	defn:
	\begin{equation*}
		\tilde{f}_{n}(x) \ = \
		\begin{cases} f_{n}(x) & x \in E \\ 0 & \trm{otw} \end{cases}
	\end{equation*}
	then $\tilde{f}(x)=\lim_{n}\tilde{f}_{n}(x)$ exists $\fa \bR$.
	note:
	\begin{enumerate}[i)]
		\item $\tilde{f}_{n}(x)$ LM since $\tilde{f}_{n}=\chi_{E}f_{n}$
		\item $\tilde{f}$ LM since its ptwise lim of LM fns
	\end{enumerate}
	thus $f=\tilde{f}$ a.e., so by lemma, $f$ is LM.
\end{pf} \

\begin{exr}
	show $E$ defn'd as above is always meas., even if the lim
	doesnt necessarily exist ae.
\end{exr}
hint: lim exists iff limsup = liminf.

\subsection{Simple functions}
throughout this section, $E\in\cl{E}$ unless stated otw.

\begin{defn}
	an LM $\phi:E\gto\bR$ is \textbf{simple} if $\abs{\trm{im}(\phi)}<\infty$.
	equivalently:
	\begin{equation*}
		\phi \ = \ \sum_{k=1}^{n}\alpha_{k}\chi_{E_{k}}
	\end{equation*}
	where $\alpha_{k}\in\bR,E_{k}\in\cl{E}$.
\end{defn} \

\begin{lm}
	sps $f:E\gto\bR$ bdd LM, $\ep>0$.
	then $\ex\phi_{\ep},\psi_{e}$ simple s.t. $\phi_{\ep}\leq f\leq\psi_{\ep}$,
	and $\abs{\psi_{e}-\phi_{e}}<\ep$ for all $x$.
\end{lm} \

\begin{pf}
	let $M>\abs{f(x)}$ for all $x\in E$, and write:
	\begin{equation*}
		(-M,M] \ = \ \bigcup_{j=1}^{n}(y_{j-1},y_{j}]
	\end{equation*}
	s.t. $\abs{y_{j}-y_{j-1}}<\ep$ for some strictly increasing $y_{j}$.
	let $E_{j}=f^{-1}((y_{j-1},y_{j}])$.
	by choice of $M$:
	\begin{equation*}
		E \ = \ \bigsqcup_{j=1}^{n}E_{j} \qquad
		\phi_{\ep} \ = \ \sum_{j=1}^{n}y_{j-1}\chi_{E_{j}} \qquad
		\psi_{\ep} \ = \ \sum_{j=1}^{n}y_{j}\chi_{E_{j}}
	\end{equation*}
\end{pf} \

\begin{lm}
	sps $f:E\gto\bar{\bR}$ LM.
	then, $\ex$ seq of simple $\phi_{n}\sto f$ ptwise s.t. $\abs{\phi_{n}}\leq\abs{f}$.

	if $f\geq0$, then furthermore $\phi_{n}$ non-neg increasing.
\end{lm} \

\begin{pf}
	general follows from non-neg case:
	note that we have $\phi_{n}\sto(f)_{+},\psi_{n}\sto(f)_{-}$ with
	$\phi_{n}\leq(f)_{+},\psi_{n}\leq(f)_{-}$.
	thus $\phi_{n}-\psi_{n}\sto(f)_{+}-(f)_{-}=f$, and:
	\begin{equation*}
		\abs{\phi_{n}-\psi_{n}}=\phi_{n}+\psi_{n}\leq(f)_{+}+(f)_{-}=\abs{f}
	\end{equation*}
	for $f\geq0$, consider $f_{n}(x)=\min(f(x),n)$ LM.
	by lemma, $\ex\phi_{n}$ simple s.t. $\phi_{n}\leq f_{n}$ and $\abs{f_{n}(x)-\phi_{n}(x)}\leq1/n$.
	let $\psi_{1}=\max(\phi_{1},0)$, and $\psi_{n}=\max(\phi_{n},\psi_{n-1})$.
	claim:
	\begin{enumerate}[i)]
		\item $\psi_{n}$ are simple:
			by ind, $\psi_{n}$ is max of two simples thus simple
			(fairly clear)
		\item $\psi_{n}$ are increasing: clear
		\item $\psi_{n}\sto f$:
			by construction + ind, $\phi_{n}\leq\psi_{n}\leq f_{n}\leq f$ for all $n$
	\end{enumerate}
\end{pf}
